\section{Introducción}
La gestión de incidencias constituye un reto significativo en instituciones públicas y privadas. Con frecuencia, los procesos de atención a solicitudes se realizan de manera manual, lo que conlleva a problemas de trazabilidad, duplicación de casos y pérdida de información crítica. Este contexto motivó el desarrollo de un Sistema de Gestión de Tickets con soporte web y móvil, cuyo propósito es centralizar la información, generar métricas confiables y ofrecer herramientas de comunicación efectivas entre técnicos, usuarios y administradores.

El Informe de Especificación de Requerimientos (SENA, 2025) sirvió como punto de partida, estableciendo las funcionalidades clave: autenticación segura, notificaciones en tiempo real, clasificación automática de tickets y un módulo de métricas. La propuesta combina tecnologías de desarrollo web y móvil de uso extendido, como React en el frontend, Laravel en el backend \cite{subecz2021} y React Native para potencial expansión multiplataforma. Este conjunto tecnológico fue seleccionado no solo por su robustez, sino también por la disponibilidad de comunidades activas y bibliografía especializada que documenta buenas prácticas y casos de éxito.

En términos metodológicos, se adoptó Scrum como marco de trabajo, el cual facilita entregas incrementales y fomenta la adaptabilidad del proyecto. La literatura destaca que las metodologías ágiles, aplicadas en conjunto con frameworks modernos, permiten reducir la deuda técnica y obtener productos más cercanos a las necesidades reales de los usuarios.

Este artículo tiene como objetivo no solo exponer el desarrollo técnico del sistema, sino también establecer un análisis comparativo con investigaciones previas. De este modo, se valida la pertinencia de las decisiones tecnológicas tomadas y se identifican oportunidades de mejora futura, tales como la integración de microservicios y la incorporación de módulos basados en inteligencia artificial.